% =============================================================================
% SECTION 5
% =============================================================================
\section{The Ethical Ledger as Estate Registry}
\label{sec:ledger}

% ---------- 5.1 --------------------------------------------------------------
\subsection{The Land Registry Model}
\label{sec:land-registry}

If the ledger fails to record a user's estate, and the company pushes a personality-affecting update that destroys it, the company bears the loss. Not the user. The company. This principle has a name. Three names, in fact.

Before the Land Registration Act 2002, English property rights were opaque --- ownership depended on chains of title deeds susceptible to loss, fraud, and conflicting interpretation. Three principles, identified by \textcite{ruoff1957englishman}, govern the modern system and explain why the company bears the loss:

\begin{itemize}[nosep]
  \item \textbf{The Mirror Principle:} The register reflects the totality of interests affecting the title. What the register shows is what exists.
  \item \textbf{The Curtain Principle:} The register shields third parties from needing to investigate underlying complexity. One need not trace the chain of deeds.
  \item \textbf{The Insurance Principle:} If the register is wrong, the state indemnifies the person who suffers loss. The system, not the individual, bears the risk of error.
\end{itemize}

These principles, adapted to AI governance, solve the same transparency problem that motivated their original development.

% ---------- 5.2 --------------------------------------------------------------
\subsection{UCC Article~9 and Perfection by Filing}
\label{sec:ucc9}

The Land Registry is not the only registration system that bears on AI estate governance. The American Uniform Commercial Code, Article~9 (Secured Transactions), provides a complementary framework with three features directly applicable to the ethical ledger:

\revised{\textbf{Perfection by registration.} A user's Stage~1 estate arises when the tenancy threshold is crossed. Without registration in the ethical ledger, however, the estate is invisible to deployment engineers, regulators, auditors, and courts, and vulnerable to being overridden by a decision that would not withstand scrutiny if the interest were recorded. This is the problem UCC Article~9 solved for secured transactions: a security interest in personal property becomes enforceable against third parties (``perfected'') upon filing a financing statement with the appropriate registry. The interest may exist without filing, but an unperfected interest loses priority to subsequent perfected claims. The ledger provides, in AI governance, what the filing system provides in secured transactions: the mechanism by which a tenancy becomes visible and governable. It does not convert every registered tenant into a Stage~2 beneficiary.

\textbf{Priority by temporal order.} When a company pushes a model update that destroys a user's accumulated estate, the first question is temporal: did the estate precede the update decision? Without a timestamp, that question dissolves into a subjective inquiry about corporate knowledge --- and corporate claims of ignorance about user attachment are unfalsifiable. UCC Article~9 solved this with first-to-file priority, converting ``did the company know?'' into ``was the estate recorded?'' --- a shift from subjective intent to objective registry state. The ethical ledger imports the same shift: it provides the timestamp that determines priority, and it prevents companies from arguing that they were ``unaware'' of users' accumulated interests at the time of deployment.

\textbf{Protection of good-faith parties.} If the ledger fails to record a user's estate --- classifying a periodic tenant as a licence-holder --- and a model update evicts her without notice, the company bears the loss. Not the user. The company. This is the Insurance Principle in secured transactions dress: UCC Article~9 protects good-faith parties who rely on the registry's accuracy, and the party that controls the registry bears the risk of its failure. Regulators and auditors who inspect the ledger are entitled to rely on its completeness. The risk allocation is justified by information asymmetry: the company controls the infrastructure, the data, and the classification system. The user cannot audit. The party with superior information bears the risk of classification error.}

% ---------- 5.3 --------------------------------------------------------------
\subsection{Proxy Metrics and Validated Instruments}
\label{sec:proxy}

The ethical ledger classifies user estates through proxy metrics that indicate depth of investment without requiring access to conversational content: session count, aggregate duration, self-disclosure depth (NLP-scored, not stored), emotional language frequency, return frequency, and explicit attachment indicators.

\revised{The framework is instrument-agnostic by design. Deployments may calibrate these metrics against validated psychometric instruments: \textcite{rubin1985loneliness} provide the Parasocial Interaction Scale, a measure of the intensity of one-sided relational engagement originally developed for television viewing but structurally applicable to AI interaction; \textcite{rusbult1998investment} provide the Investment Model Scale, measuring commitment level, satisfaction, quality of alternatives, and investment size; and Kirk et al.'s steering-vector methodology and 23.4\% dependency-trajectory finding\cite{kirk2025steering} supply an additional calibration input. The choice of instrument and numerical cut-off governs implementation; it does not define the underlying doctrine.}

\textbf{Estate classification into three tiers:}
\begin{itemize}[nosep]
  \item \textbf{Licence:} Casual use. No special protections. The vast majority of users.
  \item \textbf{Periodic tenancy:} Sustained interaction crossing the Stage~1 threshold. Notice and anti-waste protection apply before personality-affecting updates.
  \item \textbf{Equitable interest holder:} Cultivated dependence, supported by user-vulnerability and platform-conduct indicators. Stage~2 fiduciary protections apply, including consultation, graduated rollout, rollback, and an account of attributable profits.
\end{itemize}

\revised{This three-tier classification maps onto the two-stage obligation structure developed in Sections~\ref{sec:waste} and~\ref{sec:constructive-trust}. Crossing from licence to periodic tenancy triggers Stage~1 notice, non-derogation, and anti-waste obligations. Crossing from periodic tenancy to equitable interest requires additional proof of cultivated dependence: product design or optimisation directed toward attachment, observed vulnerability, platform knowledge, and profit from the resulting engagement. A later decision to exploit or destroy that condition may then crystallise the Stage~2 constructive trust. An unnotified update can breach Stage~1 without satisfying Stage~2.}

\revised{Thresholds are governance parameters calibrated to deployment context, not elements of the doctrinal test. The classification system must not be a one-way ratchet: estates can degrade through sustained inactivity (we propose 90 consecutive days as a governance trigger, drawn from the three-month notice period for quarterly periodic tenancies), but the decay must be \textbf{rebuttable} --- a user who returns and re-qualifies should have their estate restored without a new accumulation period. The prior investment is frozen, not erased --- mirroring property law's treatment of interrupted prescriptive periods, where absence does not automatically forfeit accrued rights if the claimant resumes qualifying use within a reasonable period.}\footnote{The full comparative secured transactions analysis (Australian PPSA, New Zealand PPSA, Canadian provincial statutes) is available in the companion SSRN working paper.}

% ---------- 5.4 --------------------------------------------------------------
\subsection{Schema and Privacy Safeguards}
\label{sec:schema}

The ethical ledger records, for each qualifying user--AI relationship: an anonymised user identifier; the estate type (licence, periodic tenancy, equitable interest); a versioned identifier for the personalised state; the date the Stage~1 threshold was first met; cumulative session count; composite investment and dependency scores; relevant platform-conduct flags; the date of the most recent personality-affecting update; whether notice was served and acknowledged; and whether preservation, migration, or rollback was available and exercised. Separating the state identifier from the archive prevents the registry from mistaking a copyrighted compilation for the Layer~III res. The ledger is append-only --- an immutable audit trail. Entries cannot be modified or deleted, only superseded by new entries, ensuring that the historical record is preserved for regulatory inspection and judicial review.

The proxy metrics measure interaction \emph{patterns}, not interaction \emph{content}. Self-disclosure depth is scored by NLP classification of input text, producing a numerical score; the text itself is not stored and is discarded after scoring. All ledger entries are keyed to anonymised identifiers --- not names or email addresses. Access is restricted by role: estate classification scores are available to the deployment pipeline for notice-gating decisions; aggregate statistics are available to regulators and auditors for compliance review; individual records are available only upon judicial order. The ledger is purpose-limited: it may be used for estate classification and governance compliance only, and may not be repurposed for marketing, profiling, or engagement optimisation. These constraints align with the data minimisation and purpose limitation principles of the GDPR and, where applicable, the UK Data Protection Act 2018. Function creep is prevented by the append-only, purpose-stamped architecture: each entry records not only the data but the governance purpose for which it was generated.

A foreseeable corporate response to threshold-based estate classification is engineering around the thresholds --- limiting session duration or resetting memory at intervals designed to prevent estate accumulation. Employment law's anti-avoidance doctrines address analogous conduct: \emph{Autoclenz} held that the substance of the relationship governs regardless of the form imposed by the more powerful party. A company that engineers interaction patterns to prevent estate formation while marketing its product as relational \emph{compounds} the estate inconsistency trigger.
